\documentclass{article}
\usepackage[UTF8]{ctex}
\usepackage{amsmath,amssymb,amsfonts}
\usepackage{enumitem}
\usepackage{geometry}
\geometry{a4paper, margin=1in}
% === 错题本排版优化（不参与题目标识） ===
\setlength{\parindent}{0pt}
\setlength{\parskip}{0.35em plus 0.10em minus 0.08em}
\setlength{\emergencystretch}{3em}
\setlist{itemsep=0.25em, topsep=0.40em, parsep=0pt}
\allowdisplaybreaks[2]
% === 错题本排版优化结束 ===
\begin{document}
\section*{华东师范大学2025年高等代数真题}
\section*{33．华东师范大学2025年高等代数试卷}
一．填空题。

\begin{itemize}

\subsection{题目1}
  \item[1．] 满足以下同余方程组
\[
\begin{cases}f(x) \equiv 6 & (\bmod x+1) \\ f(x) \equiv 3 x & \left(\bmod x^{2}+x+1\right) \\ f(x) \equiv(x-1)^{2} & \left(\bmod 2 x^{3}+1\right)\end{cases}
\]
且次数达到最小的多项式 \(f(x)\) 为 \(\_\_\_\_\)．

\subsection{题目2}
  \item[2．] 实系数多项式 \(f(x)=x^{5}+3 x^{4}+5 x^{3}+5 x^{2}+3 x+1\) 与其导数 \(f^{\prime}(x)\) 的最大公因式为 \(\_\_\_\_\)．
\end{itemize}

\subsection{题目3}
3．使实二次型 \(q(x, y, z)=x^{2}+y^{2}+z^{2}+t(x z-x y-y z)\) 正定的实数 \(t\) 的取值范围是 \(\_\_\_\_\)．\\
\(\_\_\_\_\)

\subsection{题目4}
4．已知方阵 \(A\) 的初等因子组为 \((\lambda+2),(\lambda+2)^{2},(\lambda-2)^{2},(\lambda-2)^{2}\) ，则其逆矩阵 \(A^{-1}\) 的极小多项式\(m_{A^{-1}}(\lambda)\) 为 \(\_\_\_\_\)．

\subsection{题目5}
5．考虑标准欧氏空间 \(\mathbb{R}^{3}\) 中的向量 \(\alpha=(1,1,1), \beta=(1,0,1), \gamma=(0,1,1)\) ，设 \(s, t\) 是使得 \(|\alpha-t \beta-s \gamma|\)达到最小的实数，那么向量 \(\alpha-t \beta\) 与 \(\alpha-s \gamma\) 的夹角余弦值为 \(\_\_\_\_\)．

\subsection{题目6}
6．设 \(A=J_{n}(\lambda) J_{n}(\mu)-J_{n}(\mu)-J_{n}(\lambda)(\lambda, \mu \in \mathbb{C})\) ，则矩阵函数 \(e^{A}\) 的行列式为 \(\_\_\_\_\)．

\begin{itemize}

\subsection{题目7}
\end{itemize}

\subsection{题目7}
原题边界“7．”：已知线性映射 \(\varphi: \mathbb{R}^{4} \rightarrow \mathbb{R}^{4}\) 满足
\[
\varphi(x, y, z, w)=(10 x-3 z+4 w, 2 x+3 y+2 w, 5 y+z+2 w,-2 x+2 y+z) .
\]
则核空间 \(\operatorname{Ker}(\varphi)\) 的维数为 \(\_\_\_\_\)．

\subsection{题目8}

\subsection{题目8}
原题边界“8．”：考虑置换 \(\sigma=(1276)(354), \tau=(1637)(24)(5)\) ，则 \(\sigma^{-1} \tau^{-1} \sigma \tau\) 等于 \(\_\_\_\_\)（写成不相交轮换积）。

\subsection{题目9}
9．实对称矩阵 \(A=J_{3}(0)+J_{3}(0)^{\mathrm{T}}\) 可通过正交相似化为对角阵 \(\_\_\_\_\)．

\subsection{题目10}
10．设 \(A=J_{2025}(0)\) ，复线性空间 \(V=\left\{X \in M_{2025}(\mathbb{C}) \mid A X=X A^{2}\right\}\) 的维数为 \(\_\_\_\_\)．

二．解答题。\\

\subsection{题目11}
1．设多项式 \(f(x)=x^{4}+2 x^{2}+a, g(x)=x^{4}-4 x^{2}+12 x-9(a \in \mathbb{C})\) 在复数域上有公共根，求 \(a\) 的值并确定 \(f(x)\) 与 \(g(x)\) 的公共根．

\subsection{题目12}
2．考虑方程组

\[
\left\{\begin{array}{l}
\lambda x_{1}+(3-\lambda) x_{2}-2 x_{3}=1 \\
(4-\lambda) x_{1}+(7-\lambda) x_{2}-6 x_{3}=3 \\
4 x_{1}+(9-\lambda) x_{2}+(\lambda-9) x_{3}=\lambda+3
\end{array}\right.
\]

问 \(\lambda\) 取何值时，方程组无解，有唯一解，有无穷多解？并在有解时求出该方程组的全部解。

\subsection{题目13}
3．设 \(A \in M_{n}(\mathbb{C})\) 满足 \(A^{n}=O\) ，但 \(A^{n-1} \neq O(n \geq 2)\) ，证明：不存在 \(n\) 阶复矩阵 \(B\) ，使得 \(A=B^{2}\) 。

\subsection{题目14}
4．设 \(A \in M_{n}(\mathbb{R})\) 是正惯性指数为1的实对称矩阵，\(V\) 是实数域上 \(n\) 维列向量空间。定义 \(V\) 上的实对称双线性型 \(f(u, v)=u^{\mathrm{T}} A v, u, v \in V\) 。设 \(u \in V\) 满足 \(f(u, u)>0\) ，证明：对任意的 \(v \in V\) ，有

\[
f(u, u) f(v, v) \leq f(u, v)^{2} .
\]

且等号成立当且仅当存在 \(r \in \mathbb{R}\) ，使得 \(f(w, v)=f(w, r u), \forall w \in V\) ．

\subsection{题目15}
5．复数域上全体二元二次型构成集合 \(V\) 。

\[
V=\left\{c_{1} x^{2}+c_{2} x y+c_{3} y^{2} \mid c_{1}, c_{2}, c_{3} \in \mathbb{C}\right\}
\]

视为以 \(x^{2}, x y, y^{2}\) 为基的复线性空间。给定矩阵 \(M=\left(\begin{array}{cc}a & b \\ c & d\end{array}\right)\) ，考虑 \(V\) 上的线性变换

\[
\varphi_{M}: V \mapsto V
\]

\[
q(x, y) \rightarrow q(a x+b y, c x+d y)
\]

（1）证明：\(\varphi_{M}\) 有一个特征值 \(\lambda_{1}=|M|\) ，且在 \(M=J_{2}(\mu)\) 的情形下求出属于 \(\lambda_{1}\) 的所有特征向量（二次型形式）。\\
（2）证明：\(\varphi_{M}\) 可对角化当且仅当 \(M\) 可对角化．
\end{document}
